%! Transformations of Functions — Unit 1, Lesson 1.4 — Warm-Up (ANSWER KEY)
\documentclass[10pt]{article}
\usepackage{precalculus-article}
\usepackage{precalculus-key}
\newcommand{\TallMath}[1]{$\displaystyle #1\rule[-1.4em]{0pt}{3.2em}$}

\begin{document}

\pageheader{Unit 1, Lesson 1.4}{Warm-Up --- Answer Key}

\vspace{0.12in}

\begin{enumerate}[itemsep=12pt]

\item From Lesson 1.1: let $f(x) = x^2$. Evaluate each one.

\begin{work}
  f(-2) &= (-2)^2 \\
        &= 4 \\
   f(1) &= (1)^2 \\
        &= 1 \\
   f(3) &= (3)^2 \\
        &= 9
\end{work}

$f(-2) = $ \ans{4} \qquad $f(1) = $ \ans{1} \qquad
$f(3) = $ \ans{9}

\item Solve each equation. Small questions today --- one of them is the whole
      explanation of a move you will meet in Part 3.

\begin{enumerate}[label=(\alph*), itemsep=5pt]
  \item $x + 3 = 0$ \qquad $x = $ \ans{$-3$}
  \item $x - 4 = 0$ \qquad $x = $ \ans{4}
  \item $2x = 6$ \qquad\; $x = $ \ans{3}
\end{enumerate}

\item From Lesson 1.1: the table gives every value of a function $p$.

\smallskip
\renewcommand{\arraystretch}{1.4}
\begin{tabular}{|c|c|c|c|c|c|}
\hline
$x$ & $-3$ & $-1$ & 0 & 1 & 3 \\
\hline
$p(x)$ & 7 & 3 & 2 & 3 & 7 \\
\hline
\end{tabular}
\smallskip

\begin{enumerate}[label=(\alph*), itemsep=5pt]
  \item Domain of $p$: \ans{$\{-3,\,-1,\,0,\,1,\,3\}$}
  \item Range of $p$: \ans{$\{2,\,3,\,7\}$}
  \item The smallest output of $p$ is \ans{2}, and it happens at
        $x = $ \ans{0}.
\end{enumerate}

\end{enumerate}

\end{document}
